% Implementierung — claude-tunnel bachelor thesis.
\chapter{Implementierung}
\label{ch:implementierung}

Der Prototyp ist in Rust als Cargo-Workspace aus fuenf Crates umgesetzt. Dieses
Kapitel beschreibt die Aufteilung (Abschnitt~\ref{sec:impl-workspace}), die
gemeinsamen Bausteine (\ref{sec:impl-common}), den Rollen-Dispatch am Edge
(\ref{sec:impl-dispatch}) sowie den Daten- und Steuerpfad
(\ref{sec:impl-datenpfad}).

\section{Workspace-Struktur}
\label{sec:impl-workspace}

\begin{table}[t]
  \centering
  \caption{Crates des Workspace und ihre Verantwortung.}
  \label{tab:crates}
  \begin{tabular}{ll}
    \toprule
    Crate & Verantwortung \\
    \midrule
    \texttt{ct-common} & Wire-Typen, Noise, PoW, Framing, Metriken \\
    \texttt{ct-edge} & providerblinder Vermittler (Rollen-Dispatch, Relay) \\
    \texttt{ct-agent} & kundenseitig, Origin-Schluessel, Serve-Pfad \\
    \texttt{ct-control-plane} & Enrollment, Registry/Rendezvous, Abrechnung \\
    \texttt{ct-client} & Tunnel-Aufbau, Betriebsarten, Bench-Harness \\
    \bottomrule
  \end{tabular}
\end{table}

Die Schichtung ist azyklisch: alle Crates haengen nur von \texttt{ct-common} ab,
nie voneinander (Testabhaengigkeiten ausgenommen). So bleibt der providerblinde
Edge unabhaengig vom kundenseitigen Agent.

\section{Gemeinsame Bausteine (ct-common)}
\label{sec:impl-common}

Der Kern liegt in \texttt{ct-common}. Der Proof-of-Work nutzt fuehrende Nullbits
eines SHA-256-Hashes; die Pruefung ist ein einzelner Hash (Listing~\ref{lst:pow}).

\begin{lstlisting}[caption={Proof-of-Work: Pruefung der Loesung (vereinfacht).},label=lst:pow]
pub fn verify(challenge: &Challenge, solution: &[u8]) -> bool {
    let mut h = Sha256::new();
    h.update(&challenge.nonce);
    h.update(solution);
    leading_zero_bits(&h.finalize()) >= challenge.difficulty
}
\end{lstlisting}

Der zentrale Datenpfad-Baustein ist \texttt{noise\_pump}: eine
transport-agnostische, bidirektionale Bruecke, die Klartext und verschluesselte
Frames in beide Richtungen konkurrent umsetzt. Der \texttt{TransportState} wird
dabei unter einem kurz gehaltenen Mutex gefuehrt, der nie ueber einen
\texttt{await}-Punkt gehalten wird. Auf diesem Baustein setzen Streaming,
Datagramm, direkter Pfad und TCP-Fallback gleichermassen auf.

\section{Rollen-Dispatch am Edge (ct-edge)}
\label{sec:impl-dispatch}

Der Edge liest je Stream ein fuehrendes Rollen-Byte und verzweigt
(Listing~\ref{lst:dispatch}). Der \texttt{'C'}-Zweig fuehrt das
PoW-Rendezvous durch und relayt danach nur noch Bytes.

\begin{lstlisting}[caption={Rollen-Dispatch (vereinfacht).},label=lst:dispatch]
match role {
    b'A' => register_agent(conn, &mut recv, state).await?,
    b'C' => rendezvous_route_relay(conn, state, challenge).await?,
    b'D' => advertise_direct(&mut recv, state).await?,
    b'P' => answer_direct_query(&mut recv, &mut send, state).await?,
    _ => return Err("unknown role".into()),
}
\end{lstlisting}

Weil Registrierung, Relay und Direktpfad-Abfrage denselben Stream-Typ nutzen,
bleibt der Edge ohne separates Steuerprotokoll erweiterbar.

\section{Daten- und Steuerpfad (agent, client, control-plane)}
\label{sec:impl-datenpfad}

Agentenseitig terminiert \texttt{serve\_noise\_stream} bzw.\
\texttt{serve\_noise\_udp} die Noise-Sitzung mit dem Origin-Schluessel und
bridget zur lokalen \gls{ac:tcp}- oder \gls{ac:udp}-Instanz; der Datenpfad ist
mit Tunnel-Metriken instrumentiert (Tunnelzahl, Bytes je Richtung,
Handshake-Latenz). Clientseitig waehlt \texttt{client\_tunnel\_auto} den direkten
Pfad, wenn ein Endpunkt advertisiert ist, und faellt sonst auf den Relay zurueck;
bei blockiertem \gls{ac:udp} traegt \texttt{client\_tunnel\_noise\_tcp} denselben
Tunnel ueber \gls{ac:tls}-ueber-\gls{ac:tcp}.

Die Steuerebene stellt Enrollment, Registry/Rendezvous und Abrechnung als
HTTP-Dienst bereit; ein \texttt{ControlPlaneClient} treibt den vollen Ablauf
(Konto eroeffnen, aufladen, Token kaufen, aufloesen) gegen den laufenden Dienst.
Die Korrektheit des Zusammenspiels ist durch eine frozen Testsuite belegt, die
den Pfad bis hin zum verschluesselten Roundtrip durch echte Komponenten
ueberprueft.
